\section{Underlying Functions and Addresses}
This section defines the cryptographic primitives and addressing scheme used throughout SHRINCS.
\subsection{Hash Function}

SHRINCS uses SHA-256 truncated to \texttt{n} bytes for all internal hash operations. 
\begin{center}
    \texttt{H(m) = SHA-256(m)[0:n]}
\end{center}

\paragraph{SHA-256 Block Precomputation} 
For SHA-256, the compression function \texttt{C(H,M)} operates on 64-byte blocks. Since \texttt{PK.seed} is 16 bytes, we pad it to 64 bytes to enable precomputation of the first block:

\begin{center}
    \texttt{P = C(IV, PK.seed || 0\string^48)}
\end{center}

Then we use the following precomputation \texttt{P} in the next compression function iteration.
Every function that uses SHA-256 with it's argument involving PK.seed will use this precomputation value. This optimization allows to reduce the number of hashing operation, which leads to faster computation, especially in the grinding parts of the SHRINCS. 

Additionally, we will denote the concatenation \texttt{PK.seed || 0\string^48} as \texttt{PK.seed*}. 

\subsection{Tweakable Hash Function}
To prevent multi-target attacks and ensure domain separation between the various components, we employ the tweakable hash function construction defined in SPHINCS+ \cite{sphincsplus_paper2019}:

\begin{center}
    \texttt{Tw(PK.seed, ADRS, m) = H(PK.seed* || ADRS || m)}
\end{center}

This usage of tweaked hashing is critical. It ensures that a hash computed at one position in the Merkle tree cannot be substituted for a hash at another position, effectively mitigating generic tree-grafting attacks.

\subsection{Pseudorandom Functions}
SHRINCS uses two pseudorandom functions for different purposes: key derivation and getting the randomizer for the signature.

For deterministic derivation of secret key material (WOTS+C chain starting values, FORS secret keys):

\begin{center}
    \texttt{PRF(PK.seed, ADRS, SK.seed) = Tw(PK.seed*, ADRS, SK.seed)}
\end{center}

The address structure \texttt{ADRS} ensures that each derived value is unique and bound to its position in the scheme.

For generating the randomizer \texttt{R} used in message hashing:

\begin{center}
    \texttt{PRF\_msg(SK.prf, PK.seed, opt\_rand, m) = HMAC-SHA-256(SK.prf, PK.seed || opt\_rand || m)[0:R\_SIZE]},
\end{center}
where \texttt{opt\_rand} is an optional randomness (\texttt{n} bytes), can be \texttt{PK.seed} for deterministic signing.
For the stateless signing path, the FORS+C grinding counter is iterated inside \texttt{PRF\_msg} to produce a fresh \texttt{R} on each attempt. The counter-augmented variant is:

\begin{center}
    \texttt{PRF\_msg\_ctr(SK.prf, PK.seed, opt\_rand, ctr, m) = HMAC-SHA-256(SK.prf, PK.seed || opt\_rand || ctr || m)[0:R\_SIZE]}
\end{center}

where \texttt{ctr} is a counter incremented during FORS+C grinding (see Section~9.4). The verifier never sees \texttt{ctr}, it only receives the final \texttt{R} that satisfied the grinding condition. This design keeps the grinding work inside the signer's PRF evaluation while allowing the verifier to check the result using only \texttt{R}  and the message hash.

\subsection{Message Hash Function}
SHRINCS uses different hashing strategies for WOTS+C and FORS+C, optimized for their respective use cases.

\subsubsection{WOTS+C: Two-Phase Hashing}
WOTS+C uses a two-phase approach for message hashing and grinding. This design is necessary because:
\begin{itemize}
    \item User messages need randomized hashing (phase 1 + phase 2)
    \item Internal XMSS signatures sign tree roots that are already hashes (phase 2 only)
\end{itemize}

\paragraph{Phase 1. Message Digest.} Computes a randomized hash of the user message:
\begin{center}
    \texttt{H\_msg(ADRS, R, PK.seed, PK.root, m) = SHA-256(PK.seed* || ADRS || R || PK.root || m)[0:n]}
\end{center}
This is computed once per signature, producing an n-byte digest.

\paragraph{Phase 2. Grinding Hash.} Takes the digest and a counter:
\begin{center}
    \texttt{H\_grind(ADRS, PK.seed, digest, ctr) = SHA-256(PK.seed* || ADRS || digest || ctr)[0:n]}
\end{center}

\subsubsection{FORS+C: Single-Phase Hashing}
FORS+C always signs user messages (never tree roots), so it uses a simpler single-phase approach. The grinding counter is not included in the message hash. Instead, it is iterated inside \texttt{PRF\_msg\_ctr} (Section~5.3) to produce a fresh \texttt{R} on each grinding attempt. The verifier only sees the final \texttt{R}:
\begin{center}
    \texttt{H\_msg\_fors(ADRS, R, PK.seed, PK.root, m) = SHA-256(PK.seed* || ADRS || R || PK.root || m)[0:roundup((k*a + hsl) / 8)]}
\end{center}

The output length is roundup((k*a + hsl) / 8) = 20 bytes for both parameter sets, encoding both the \texttt{k} FORS leaf indices and the \texttt{hsl} hypertree path indices in a unified digest (see Section~10.5).

\subsection{Tweak Structure and Domain Separation}
Tweaks are structured as follows to ensure uniqueness across all scheme components:

\begin{center}
\texttt{ADRS = layer || tree\_address || type || words}
\end{center}
where: 
\begin{itemize}
    \item[] \texttt{layer}: (4 bytes) layer in hypertree (\texttt{0} = bottom, \texttt{d-1} = max)
    \item[] \texttt{tree\_address}: (12 bytes) tree address within layer
    \item[] \texttt{type}: (4 bytes) address type
    \item[] \texttt{words}: (12 bytes) type-dependent 12 bytes
\end{itemize}
A total \texttt{ADRS} size is 32 bytes.

\subsubsection{Address Types}
\begin{verbatim}
0x00: SF_WOTS_HASH:     Stateful WOTS+C chain hashing
0x01: SF_WOTS_PK:       Stateful WOTS+C public key compression
0x02: SF_TREE:          Stateful tree internal nodes
0x03: SF_WOTS_GRIND:    Stateful WOTS+C grinding
0x04: SF_H_MSG:         Stateful message hash
0x05: SF_WOTS_PRF:      Stateful PRF for WOTS+C key derivation
0x06: FORS_HASH:        FORS leaf hashing
0x07: FORS_TREE:        FORS tree internal nodes
0x08: FORS_PK:          FORS roots compression
0x09: FORS_PRF:         PRF for FORS secret key derivation
0x0A: SL_WOTS_HASH:     Stateless WOTS+C chain hashing
0x0B: SL_WOTS_PK:       Stateless WOTS+C public key compression
0x0C: SL_TREE:          Stateless tree internal nodes
0x0D: SL_WOTS_GRIND:    Stateless WOTS+C grinding
0x0E: SL_H_MSG:         Stateless message hash
0x0F: SL_WOTS_PRF:      Stateless PRF for WOTS+C key derivation
0x10: ROOT:             Root public key computation
\end{verbatim}
Summarizing, \texttt{(0x00, 0x01, 0x02, 0x03, 0x04, 0x05)} are stateful only types, \texttt{(0x06, 0x07, 0x08, 0x09, 0x0A, 0x0B, 0x0C, 0x0D, 0x0E, 0x0F)} are stateless only and \texttt{(0x10)} is shared type. 

\subsubsection{Address Manipulation Functions}
\begin{verbatim}
new_ADRS():
     return toByte(0, 32)

setLayerAddress(ADRS, layer):
    ADRS[0:4] ← toByte(layer, 4)

setTreeAddress(ADRS, tree_addr):
    ADRS[4:16] ← toByte(tree_addr, 12)

setTypeAndClear(ADRS, type):
    ADRS[16:20] ← toByte(type, 4)
    ADRS[20:32] ← toByte(0, 12)

setKeyPairAddress(ADRS, keypair):
    ADRS[20:24] ← toByte(keypair, 4)

setChainAddress(ADRS, chain):
    ADRS[24:28] ← toByte(chain, 4)

setHashAddress(ADRS, hash):
    ADRS[28:32] ← toByte(hash, 4)

setTreeHeight(ADRS, height):
    ADRS[24:28] ← toByte(height, 4)

setTreeIndex(ADRS, index):
    ADRS[28:32] ← toByte(index, 4)
\end{verbatim}

Note that \texttt{setTreeHeight} and \texttt{setChainAddress} both write to bytes \texttt{[24:28]}. These operations are context-dependent so there mustn't be a conflict.

\subsubsection{Type-Dependent Words}
The list of type-dependent words for address tweaking is the following:
\begin{table}[h!]
\centering
\begin{tabular}{c c c c c}
\hline
\textbf{Type} & \textbf{Name} & \textbf{word1 [20:24]} & \textbf{word2 [24:28]} & \textbf{word3 [28:32]} \\
\hline

0x00 & \texttt{SF\_WOTS\_HASH}   & \texttt{keypair}   & \texttt{chain}   & \texttt{hash} \\
0x01 & \texttt{SF\_WOTS\_PK}     & \texttt{keypair}   & \texttt{0}       & \texttt{0} \\
0x02 & \texttt{SF\_TREE}         & \texttt{keypair}   & \texttt{height}  & \texttt{index} \\
0x03 & \texttt{SF\_WOTS\_GRIND}  & \texttt{keypair}   & \texttt{0}       & \texttt{0} \\
0x04 & \texttt{SF\_H\_MSG}       & \texttt{0}         & \texttt{0}       & \texttt{0} \\
0x05 & \texttt{SF\_WOTS\_PRF}    & \texttt{keypair}   & \texttt{chain}   & \texttt{0} \\
0x06 & \texttt{FORS\_HASH}       & \texttt{tree\_idx} & \texttt{0}       & \texttt{leaf\_idx} \\
0x07 & \texttt{FORS\_TREE}       & \texttt{tree\_idx} & \texttt{height}  & \texttt{index} \\
0x08 & \texttt{FORS\_PK}         & \texttt{0}         & \texttt{0}       & \texttt{0} \\
0x09 & \texttt{FORS\_PRF}        & \texttt{tree\_idx} & \texttt{0}       & \texttt{leaf\_idx} \\
0x0A & \texttt{SL\_WOTS\_HASH}   & \texttt{keypair}   & \texttt{chain}   & \texttt{hash} \\
0x0B & \texttt{SL\_WOTS\_PK}     & \texttt{keypair}   & \texttt{0}       & \texttt{0} \\
0x0C & \texttt{SL\_TREE}         & \texttt{keypair}   & \texttt{height}  & \texttt{index} \\
0x0D & \texttt{SL\_WOTS\_GRIN}   & \texttt{keypair}   & \texttt{0}       & \texttt{0} \\
0x0E & \texttt{SL\_H\_MSG}       & \texttt{0}         & \texttt{0}       & \texttt{0} \\
0x0F & \texttt{SL\_WOTS\_PRF}    & \texttt{keypair}   & \texttt{chain}   & \texttt{0} \\
0x10 & \texttt{ROOT}             & \texttt{0}         & \texttt{0}       & \texttt{0} \\

\hline
\end{tabular}
\caption{Type-dependent interpretation of ADRS words \texttt{(bytes [20:32])}}
\end{table}